\section{Hodge conservation after one stabilization}\label{sec:hodge-conservation}

The numerical argument selected whole quantum primary factors whose
contributions vanish on surface centers. We now retain the odd Hodge
representation of those factors. The new point is to restrict the quantum
parameters while keeping the entire cohomology fiber.

We prove Theorem~\ref{thm:hodge-conservation} for the nine families
$\mathcal F$ listed in Proposition~\ref{prop:fano-endpoints}.

Part~\ref{part:numerical} supplies four ingredients used here: regular separation
of whole primary connection blocks; cyclic persistence in all even bulk
directions and rank-two residue transport; the low-dimensional vanishing
calculations; and the nine-family endpoint computations and specialized
$\PP^1$ lemma (Lemmas~\ref{lem:cyclic-primary-persistence},
\ref{lem:triple-primary-persistence} and~\ref{lem:special-projective-line},
Propositions~\ref{prop:rank2-rigidity}, \ref{prop:atomic-lowdim}
and~\ref{prop:fano-endpoints}). We reuse these proofs on the bases specified below. The
fixed-base comparison and the final recovery of a rational Hodge structure
are proved in this section. No assertion of the numerical classification
requires this theorem.

\subsection{A common group and the retained fiber}

For a polarizable rational Hodge structure, the Hodge group is the rational
algebraic group generated by the circle action that acts on $H^{p,q}$ as
$\lambda^{p-q}$. Tate twists do not change this action. Iritani recalls
its universal form in \cite[Section~3]{IritaniHodge}: the category of polarizable
rational Hodge structures modulo Tate twists has a rational fiber functor,
and its tensor automorphisms form the universal Hodge group.

Fix the finite collection of varieties and centers in a projective weak
factorization of the proposed birational map, together with the auxiliary
surface models used below. Use one reductive rational
quotient $G$ of this universal group through which all their cohomology
representations factor. The image on each cohomology space is its Hodge
group. Thus the choices for different varieties are compatible. Tate
suspensions in the blowup formula act trivially on these representations.

For each variety $T$, restrict the bulk parameter to
\[
 H^*(T,\Q)^G\otimes_\Q\C,
 \qquad\text{while the fiber remains }H^*(T,\Q)\otimes_\Q\C.
\]
The fixed parameters are even: the element $-1$ of the Hodge circle acts
by cohomological parity. The unit and algebraic divisor classes are fixed.
No assertion that every Hodge class is algebraic is needed. Novikov
parameters, $z$ and the coefficient ramification parameters have trivial
$G$-action and even parity.

The quantum product is equivariant under simultaneous action on parameters
and fibers \cite[Lemma~7]{IritaniHodge}. On this fixed base its Euler operator
therefore commutes with $G$. Its whole primary projectors commute with $G$
as well. All ranks below are full even fiber ranks, never dimensions of
$G$-invariant subspaces.

After a splitting extension of the generic coefficient field, record the
odd part of each primary factor by its multiplicities in the irreducible
representations of $G_\C$. These give a class in the representation group
$R(G_\C)$, a free abelian group on simple representations. Further
algebraically closed coefficient extension does not change these
multiplicities.

\subsection{Faithfulness on the fixed base}

The comparison must recover the intrinsic generic center connection, even
when its raw Novikov map identifies curve classes. We check faithfulness
on the fixed base directly.

\begin{lemma}[Fixed-base comparison]
\label{lem:hodge-fixed-comparison}\coverage{absent}%
For a blowup $\Bl_ZT\to T$ with smooth connected center of codimension
$c\geq2$, the actual comparison on Hodge-fixed bulk bases preserves whole
primary factors, their full $G$-representations and their original
$z$-lattices. Each center occurrence is a faithful extension of its
intrinsic reduced fixed-base connection. Eligible rank-two factors retain
their canonical modified residue discriminant.
\end{lemma}

\begin{proof}
Iritani's blowup isomorphism intertwines the quantum connections over the
specified graded completed $\C[z]$-modules
\cite[Theorem~5.18 and Section~5.8.2]{IritaniBlowup}. Its coordinate maps and full-fiber
isomorphism are equivariant, and their initial shifts are fixed
\cite[Proposition~8 and Lemma~10]{IritaniHodge}. The joint coordinate map and its
inverse restrict to inverse maps of fixed loci. In the external direct
sum the fixed coordinates of the ambient copy and each center copy are
independent; in particular their unit and divisor coordinates survive.

Write $d$ for an effective integral numerical curve class of $Z$. The
reduced source retains the combined generator
$X_d=Q_Z^d\exp(\sigma^{(2)}\cdot d)$, with no separate source divisor
variables. If $u$ denotes fixed higher-degree coordinates and $v$ the unit
coordinate, its ring is
\[
 R_Z^G=\C[[X_d,u]]_{\mathrm{gr}}[v].
\]
Here graded completion retains only finitely many total homogeneous
degrees in each element, and unit dependence is polynomial. At a fixed
curve class the coefficient is consequently polynomial in $u,v$: every
$u$ has strictly negative degree. These are precisely the reduced-domain
conventions of the numerical comparison proof.

For the $j$th occurrence, with independent fixed divisor coordinate
$s_j^{(2)}$ and fixed initial shift $\varsigma_j^\circ$, the map is
\[
 X_d\longmapsto Q^{i_*d}
 q^{-c_1(N_{Z/T})\cdot d/(c-1)}
 \exp\bigl((\varsigma_j^{\circ,(2)}+s_j^{(2)})\cdot d\bigr).
\]
The other coordinates undergo the corresponding fixed shifts. Existence
of this reduced substitution uses the string and divisor equations; it
does not authorize translation of unrestricted formal power series.

To prove injection, take the least nonzero degree with respect to the
restriction of an ambient ample divisor. Bounded numerical effective-degree
slices are finite. At this least degree, positive-degree corrections to
the initial shifts cannot contribute; translation of the polynomial
coefficients by their degree-zero shifts is injective. Group terms with
the same untagged target monomial. A kernel element would give a nonzero
finite relation
\[
 \sum_d a_d\exp\!\left(\sum_k(D_k\cdot d)x_k\right)=0,
\]
where the $a_d$ are independent of the divisor coordinates $x_k$.
Choose integral algebraic divisors $D_k$ separating numerical curve
classes. Their coordinates remain available on the fixed base. Restrict
to an integral line in these coordinates separating the finitely many
exponent vectors. Successive derivatives at zero give an invertible
Vandermonde matrix, forcing all $a_d=0$. This is a contradiction.

The original comparison and its inverse are regular at $z=0$ in the
coefficientwise realization of the graded completion. Leading-term
intertwining therefore carries $N$ and $\im N$ to their counterparts,
and hence carries the canonical modification
\[
 E^\sharp=\{s\in E:s\bmod z\in\im N\}
\]
to the corresponding modification. In modified frames the regular inverse
still exists, so the residues are conjugate. This applies also to resonant
rank-two factors. Finally, independent occurrence-unit shifts separate
their spectra generically: the resultant of two scalar-translated
characteristic polynomials is a nonzero polynomial in the shift difference.
Thus distinct occurrences do not merge in the comparison.
\uses{lem:faithful-center-base-change,prop:rank2-rigidity}%
\imports{IritaniBlowup:thm-5.18,IritaniHodge:equivariance}%
\end{proof}

The same argument applies to the ambient copy and componentwise to a
disconnected center. Each operation uses its own comparison field; the
resulting equalities of representation multiplicities telescope without
composing different Laurent expansions along the factorization.

\subsection{Selection and vanishing}

Let $E$ run over the whole primary factors of the full connection on the
fixed base. Retain $E$ if either
\[
 \rk E^\ev=2,\quad N_E^2=0\ne N_E,\quad\delta^\sharp(E)\ne0,
 \qquad\text{or}\qquad \rk E^\ev=3.
\]
Here $N_E$ is the centered leading operator on the even part, and
$\delta^\sharp$ is the canonical rank-two residue discriminant. Define
\[
 \mathcal H_G(T)=\sum_{E\text{ retained}}[E^\odd]\in R(G_\C).
\]
There is no factor of one half. Selection precedes extraction of the odd
representation. The sum has no scalar eigenvalue or discriminant labels;
permuting splitting branches does not change it. We do not require each
individually labelled factor to descend to $\Q$.

\begin{proposition}[Operation and vanishing]
\label{prop:hodge-operations}\coverage{absent}%
The selected class satisfies
\[
 \mathcal H_G(\Bl_ZT)=\mathcal H_G(T)+(c-1)\mathcal H_G(Z).
\]
It is zero for smooth projective varieties of dimension at most two and
is unchanged under birational maps of smooth projective fourfolds.
\end{proposition}

\begin{proof}
The fixed-base comparison preserves the selection conditions and the odd
representations, with the $c-1$ center Tate suspensions forgotten by $G$.
Generic separation of occurrences makes their contributions add.

The low-dimensional calculations of Propositions~\ref{prop:atomic-lowdim}
and Lemma~\ref{lem:odd-selector} apply to these bases as
follows. A point has only an even-rank-one factor. A curve has simple even
factors, zero centered operator, or a rank-two factor with discriminant
zero. For a minimal surface with nef canonical class, the centered Euler
operator strictly raises ordinary degree, also on the fixed base. Its
only primary factor has full even rank $b_2+2$. The only possible selected
rank is three, when $b_2=1$. In that case a nonzero holomorphic one-form
$\alpha$ would give a nonzero class $[i\alpha\wedge\bar\alpha]$ of square
zero, impossible in the one-dimensional second cohomology generated by
an ample class. Thus $H^1$ and $H^3$ vanish, and its odd contribution is zero.

On $C\times\PP^1$ with $g(C)\geq1$, the full-even-bulk calculation in
Lemma~\ref{lem:ruled-product} has two even-rank-two factors, each with centered operator
$(2-2g(C))h$. For genus one it is identically zero; for larger genus the
modified residue discriminant is zero. All even classes here are Tate,
so this is also the fixed-base calculation. The horizontal Novikov
parameter remains an independent spectator. For $C=\PP^1$ the generic
factors are simple. Every ruled surface is birational to such a product;
surface factorization uses point centers only, whose contributions have
already been shown to vanish. The $\PP^2$ calculation and point blowups
complete the surface classification argument. No fourfold invariance was
used to prove these vanishings.

Projective weak factorization of fourfolds now uses only smooth projective
centers of dimension at most two. Every correction is zero, proving the
last assertion.
\uses{lem:hodge-fixed-comparison,prop:atomic-lowdim,lem:ruled-product,lem:odd-selector}%
\imports{AKMW:weak-factorization,BeauvilleSurfaces:minimal-surface-classification}%
\end{proof}

\subsection{Endpoint recovery and the proof of the theorem}

We first explain why continuation preserves the information just retained.
On a connected fixed-base germ an equivariant projector has constant
representation multiplicities. Indeed, write the full fiber as a sum
$\bigoplus_i V_i\otimes M_i$ of irreducible $G_\C$-representations and
multiplicity spaces. The projector acts by the identity on $V_i$ and an
idempotent on $M_i$. Its image has locally constant rank, hence constant
rank on the germ. This argument concerns a separated continued projector;
it does not continue a decomposition through a collision of eigenvalues.

For $X\in\mathcal F$, all even cohomology is Tate and
$H^\odd(X)=H^3(X)$. The nine endpoint calculations in Proposition~\ref{prop:fano-endpoints} give one
repeated even primary factor: rank two with nonzero discriminant in
index-two degrees $1,2,3$ and index-one genera $6,8$, or rank three in
index-one genera $2,3,4,5$. All complements have even rank one.

An even-rank-one factor of a Frobenius superalgebra has no odd part.
Writing its even unit as $e$, odd elements satisfy $ab=ce$; supercommutativity
gives $(ab)^2=0$, hence $c=0$. Nondegeneracy of the odd Frobenius pairing
then forces the odd part to vanish. Consequently the entire $H^3(X)$
lies in the unique retained repeated factor, and
\[
 \mathcal H_G(X)=[H^3(X,\Q)_\C].
\]

At the small product point, the genus-zero product formula identifies the
full connection and its original lattice for $X\times\PP^1$ with the
tensor product \cite[formula~(1)]{Behrend}. The independent fiber Novikov
parameter separates the two $\PP^1$ eigenvalues and every shifted
$X$ cluster. The specialized $\PP^1$ Lemma~\ref{lem:special-projective-line} continues these
clusters in all even bulk directions, including mixed ones: cyclicity and
even rank persist, and the rank-two discriminant is unchanged. Tensoring
with the regularly separated rank-one factor changes the residue only
by a scalar. On the fixed base the equivariant-projector argument above
preserves the odd multiplicities. The Tate-shifted second copy has the
same $G$-representation. Therefore
\[
 \mathcal H_G(X\times\PP^1)=2[H^3(X,\Q)_\C].
\]
This uses neither a big quantum tensor formula at mixed bulk nor a general
projective-bundle comparison theorem.

\Needspace{15\baselineskip}
\begin{proof}[Proof of Theorem~\ref{thm:hodge-conservation}]
Choose the common group $G$ for a weak factorization between the two
products. Operation and vanishing, followed by endpoint recovery, give
\[
 2[H^3(X,\Q)_\C]=2[H^3(Y,\Q)_\C]\quad\text{in }R(G_\C).
\]
Semisimplicity makes this group free on simple classes. Cancel two and
obtain an isomorphism of the complexified $G$-representations.

Put $V=H^3(X,\Q)$ and $W=H^3(Y,\Q)$. Both have pure weight three.
On their Hodge summands the circle records $p-q$, and the common weight
records $p+q=3$; together these determine $(p,q)$. Thus a rational
$G$-equivariant map between $V$ and $W$ is precisely a rational Hodge
morphism. Equivalently, no nonzero Tate twist can identify weight-three
constituents here.

Finally, equivariant Hom commutes with scalar extension, so
\[
 \Hom_G(V,W)\otimes_\Q\C
   =\Hom_{G_\C}(V_\C,W_\C).
\]
The determinant polynomial on the finite-dimensional rational vector space
$\Hom_G(V,W)$ is nonzero because the complexified space contains an
isomorphism. Over the infinite field $\Q$, a nonzero polynomial has a
rational point where it does not vanish. That point gives the required
rational Hodge isomorphism.
\proves{thm:hodge-conservation}%
\uses{prop:hodge-operations,prop:fano-endpoints,lem:special-projective-line,lem:odd-selector}%
\imports{IritaniHodge:equivariance,AKMW:weak-factorization,Behrend:product}%
\end{proof}

The conclusion retains neither a distinguished integral lattice nor a
principal polarization. Those structures are separate inputs in any
subsequent Torelli or arithmetic application.

